% Generated draft: requires mathematical and source review.
\documentclass[11pt]{article}
\usepackage[margin=25mm]{geometry}
\usepackage{fontspec}
\setmainfont{Latin Modern Roman}
\usepackage{amsmath,amssymb,mathtools,unicode-math}
\setmathfont{Latin Modern Math}
\usepackage{enumitem,needspace,xurl,hyperref}
\setlength{\parindent}{0pt}
\setlength{\parskip}{5pt}
\setlength{\emergencystretch}{3em}
\newcommand{\N}{\mathbb{N}}
\newcommand{\Z}{\mathbb{Z}}
\newcommand{\Q}{\mathbb{Q}}
\newcommand{\R}{\mathbb{R}}
\newcommand{\C}{\mathbb{C}}
\newcommand{\F}{\mathbb{F}}
\newcommand{\A}{\mathbb{A}}
\newcommand{\PP}{\mathbb{P}}
\newcommand{\E}{\mathbb{E}}
\newcommand{\eps}{\varepsilon}
\newcommand{\dd}{\,\mathrm{d}}
\begin{document}
\section*{UnsolvedMath ID: 800879}
\textbf{Source problem code:} catalogue entry
\subsection{Definitions and mathematical statement}
The supplied excerpt does not explicitly specify the ambient geometry or define its function classes. Under the standard compact Kähler interpretation, let \(X\) be a compact Kähler manifold of complex dimension \(n\geq 1\), let \(\omega\) be a Kähler form, and put \(V=\int_X\omega^n\). These ambient hypotheses and the definitions below require confirmation against the complete source.

Fix the convention \(dd^c=\frac{i}{\pi}\partial\bar\partial\). Write \(\operatorname{PSH}(X,\omega)\) for the upper semicontinuous, locally integrable functions \(\varphi:X\to[-\infty,\infty)\) satisfying \(\omega+dd^c\varphi\geq 0\) in the sense of currents. For bounded such functions, the measure \((\omega+dd^c\varphi)^n\) is understood in the Bedford--Taylor sense.

The global domain of definition \(DMA(X,\omega)\) consists of those \(\varphi\in\operatorname{PSH}(X,\omega)\) for which there exists a positive Borel measure \(MA(\varphi)\) such that, for every sequence \(\varphi_j\in\operatorname{PSH}(X,\omega)\cap L^\infty(X)\) decreasing pointwise to \(\varphi\),
\[
(\omega+dd^c\varphi_j)^n\rightharpoonup MA(\varphi).
\]
In particular, the limiting measure must be independent of the approximating sequence.

For \(\varphi\in\operatorname{PSH}(X,\omega)\), set \(\varphi^{(j)}=\max\{\varphi,-j\}\). The full non-pluripolar Monge--Ampère mass class is
\[
\mathcal E(X,\omega)=\left\{\varphi\in\operatorname{PSH}(X,\omega):\lim_{j\to\infty}\int_{\{\varphi>-j\}}(\omega+dd^c\varphi^{(j)})^n=V\right\}.
\]

The question attributed in the supplied source to Coman--Guedj--Zeriahi asks whether, for every \(\varphi\in DMA(X,\omega)\),
\[
MA(\varphi)\bigl(\{x\in X:\varphi(x)=-\infty\}\bigr)=0
\quad\Longrightarrow\quad
\varphi\in\mathcal E(X,\omega).
\]
Here the hypothesis concerns only the pole set of \(\varphi\); it does not assert that \(MA(\varphi)\) vanishes on every pluripolar set. The supplied registry labels the question open; this draft does not independently verify that status.
\subsection{Short English statement}
In the standard compact Kähler setting, suppose the Monge--Ampère measure of a function in the global domain of definition gives no mass to the set where that function equals \(-\infty\). Must the function have full non-pluripolar Monge--Ampère mass? The excerpt leaves the ambient hypotheses and class definitions implicit, so the stated interpretation requires source confirmation.
\subsection{Sources}
\begin{itemize}
\item Supplied local UnsolvedMath registry record: permanent ID \(800879\), exact source code \texttt{catalogue entry}, title ``Open Problems in Pluripotential Theory --- Question 10.'' The registry attributes the problem to Dinew, Guedj, and Zeriahi, gives proposed year 2015, and labels it open. These are supplied metadata, not independently verified conclusions.
\item Dinew, Guedj, and Zeriahi, \emph{Open Problems in Pluripotential Theory} (2015), Question 10: \url{https://arxiv.org/abs/1511.00705}. The supplied provenance records access on 2026-10-16 and extraction by \texttt{arxiv\_latex}; neither access nor extraction was independently performed for this draft.
\item The supplied question carries the attribution ``Coman--Guedj--Zeriahi'' and cites the reference label [CGZ08]. The excerpt describes the investigation in [CGZ08] as having limited success; this is the supplied survey's assessment. No full bibliographic entry or separate URL for [CGZ08] is supplied.
\end{itemize}
\end{document}
